% Challenge report - Deep Learning (SRH). IEEE conference format, 2-4 pages incl. figures and tables
% (references do not count). Compile: pdflatex report, bibtex report, pdflatex report, pdflatex report
% - or upload the whole folder to Overleaf (New Project -> Upload Project).
\documentclass[conference]{IEEEtran}
\usepackage{graphicx}
\usepackage{booktabs}
\usepackage{amsmath}
\usepackage{url}

\begin{document}

\title{Your Title: What You Did on DermaMNIST}

\author{\IEEEauthorblockN{First Author}
\IEEEauthorblockA{SRH University\\ first.author@example.com}
\and
\IEEEauthorblockN{Second Author}
\IEEEauthorblockA{SRH University\\ second.author@example.com}}

\maketitle

\begin{abstract}
3--5 sentences: the task, your baseline, your improvements, your main result on the test set
(balanced accuracy, mean $\pm$ std over seeds), and the one thing you learned.
\end{abstract}

\section{Introduction}
Why is this task relevant? What is hard about it (7 classes, strong imbalance)?
What did you try, and what is your main finding? (about 1/3 column)

\section{Data}
DermaMNIST~\cite{medmnistv2}: number of images per split, the 7 classes, the class counts (a small table or
bar plot), example images. How did you preprocess the images (normalisation, augmentation)?

\section{Methods}
\subsection{Baseline (A)}
The CNN of the starter notebook: layers, number of parameters, loss, optimiser, learning rate, epochs.

\subsection{Improvements (B)}
Describe every change and \emph{why} you made it. At least one must be motivated by a paper or another good
source --- cite it~\cite{he2016resnet}.

\subsection{Evaluation}
Main metric: balanced accuracy. Also macro-F1 and the sensitivity per class. Models and settings were chosen
on the validation set; every variant was trained with at least 3 seeds. Differences to the baseline:
95\,\% bootstrap interval of the difference on the validation set.

\section{Results}
Table~\ref{tab:results} shows mean $\pm$ std over seeds. Change one thing at a time.

\begin{table}[t]
\caption{Balanced accuracy (mean $\pm$ std over 3 seeds)}
\label{tab:results}
\centering
\begin{tabular}{lccc}
\toprule
variant & val & test & 95\,\% interval vs.\ A (val) \\
\midrule
A: baseline & 0.46 $\pm$ 0.02 & -- & -- \\
B1: class weights & 0.53 $\pm$ 0.01 & -- & [0.02, 0.13] \\
B2: \dots & & & \\
\bottomrule
\end{tabular}
\end{table}

Add one key figure, e.g. the confusion matrix of your final model or the learning curves:
% \begin{figure}[t]
%   \centering
%   \includegraphics[width=\columnwidth]{confusion_matrix.pdf}
%   \caption{Confusion matrix of the final model on the test set.}
%   \label{fig:cm}
% \end{figure}

\section{Discussion}
What worked, what did not, and why? Which classes are still hard? What would you try next?
Limitations (small images, one dataset, \dots).

\section{Conclusion}
2--3 sentences.

\section*{AI use statement}
Which AI tools did you use, and for what (e.g.\ ``ChatGPT to explain an error message'',
``Copilot for plotting code'')? If none: ``No AI tools were used.''

\section*{Contribution statement}
Who did what (2--3 sentences), e.g.\ ``A.\ wrote the data pipeline and the augmentation experiments;
B.\ implemented the class weights and the evaluation; both wrote the report and the poster.''

\bibliographystyle{IEEEtran}
\bibliography{references}

\end{document}
